\chapter{Research brief and navigation}\label{ch:orientation}
\section{The exact problem}
For positive integers define
\[
 h(k)=k\sigma(k),\qquad \sigma(k)=\sum_{d\mid k}d,\qquad
 f(N)=\#\{k\ge1:h(k)=N\}.
\]
The primary question, recorded as Erd\H{o}s Problem 1060, is
\begin{equation}\label{eq:EP}
 \boxed{\forall\varepsilon>0\ \exists N_\varepsilon\ \forall N\ge N_\varepsilon:
 f(N)\le\exp\!\left(\varepsilon\frac{\log N}{\log\log N}\right).}
\end{equation}
The stronger, optional question asks whether $f(N)\le(\log N)^{O(1)}$.
All logarithms in this dossier are natural. We write $L=\log N$ and
$w=\log L=\log\log N$; asymptotic formulas involving further logarithms are
understood only above a sufficiently large absolute threshold.

\begin{statusbox}{Proof status at handoff}
The unrestricted estimate \eqref{eq:EP} has not been proved or refuted in this
investigation. Several full proofs of partial results, exact finite certificates,
and conditional reductions have been obtained. None of the unproved uniform
counting estimates below may be inserted as a lemma in a claimed solution.
There is no justified probability estimate that a completion is imminent.
\end{statusbox}

The equivalent target is
\[
 \log\max\{1,f(N)\}=o(L/w).
\]
An average estimate, a density-one result, $f(N)=N^{o(1)}$, or
$\log f(N)\le(C+o(1))L/w$ with a fixed $C>0$ does not prove it.
The maximum must include exceptional smooth targets and targets with very large
small-prime valuations. A bound on input exponents is not a bound on target
exponents: many local divisor sums can contribute to the same target prime.

\section{How to use this dossier}
This is an edited research handoff, not a transcript and not a submission-ready
resolution. Full working proofs have been retypeset by topic; repeated historical
claims of being close to a solution have been removed. Superseded bounds are
retained only when their mechanism is reusable. Original notes, archived code
packets, and check outputs are preserved in the accompanying handoff archive.

A productive first reading is Chapters~\ref{ch:foundations}, \ref{ch:reductions},
\ref{ch:twoprime}, \ref{ch:primitive}, \ref{ch:exchanges}, and
\ref{ch:next}. The strongest general numerical bound is in
Chapter~\ref{ch:chain}. Before adopting any graph or linear-relaxation route,
read Chapters~\ref{ch:ratio}, \ref{ch:integer}, and \ref{ch:bh}.

The working proofs have not received an independent human referee review or
formal proof-assistant verification. Compilation involved source reconciliation,
selected external-source checks, and reruns of the verifiers listed in
Chapter~\ref{ch:repro}. These are not a substitute for independently auditing a
proof before using it. No claim of novelty relative to the complete literature
has been established.

\section{Principal results and their limitations}
\small
\begin{longtable}{@{}p{0.17\textwidth}p{0.57\textwidth}p{0.18\textwidth}@{}}
\toprule
Result & Exact scope & Location\\\midrule\endhead
Powerful core & $f(N)\le\prod_{p\mid N}v_p(N)$. Known direct majorant; maximal-order constant $\log3/3$. & Ch.~\ref{ch:foundations}\\
Prime-index encoding & Fixing $v_\ell(k)$ and $(v_p(k)+1)_{\ell'}$ for $p\ne\ell$ determines a preimage. & Ch.~\ref{ch:weighted}\\
Weighted bound & General constant $C_0=\tfrac12\log3-\tfrac13\log2\approx0.318257$. Numerically superseded; local entropy mechanism retained. & Ch.~\ref{ch:weighted}\\
Chain packing & Strongest unrestricted bound in this record: $C_1=\tfrac12\log(1+1/\sqrt2)\approx0.267400$. Still positive. & Ch.~\ref{ch:chain}\\
Cubefree bound & $\log f_2(N)\le(\tfrac12\log\varphi+o(1))L/w$, $\varphi=(1+\sqrt5)/2$. & Ch.~\ref{ch:chain}\\
Compatible families & $\log|\mathcal F|=O_E(w\log w)$ with cap $E$; $O(w^{3/2})$ without a cap. Compatibility is essential. & Ch.~\ref{ch:sparse}\\
Restricted targets & Full $f(N)$ has zero constant on sixth-power-free targets, and on $N=2^a m$ with $m$ odd fourth-power-free and arbitrary $a$. & Ch.~\ref{ch:sparse}\\
Two-prime rigidity & For a fixed pair $p\ne q$, all $h(p^a q^b)$ are distinct, with unrestricted exponents. Different preimages differ on at least three bases. & Ch.~\ref{ch:twoprime}\\
Primitive-divisor lists & $f(N)\le\omega(N)^{\omega(N)}$ and each input base has at most $\omega(N)-1$ admissible exponents $>1$. & Ch.~\ref{ch:primitive}\\
Exact integer models & Complete finite domains, valuation rows, sound propagation, rational certificates, rank bounds and exact branching costs. Uniform costs unproved. & Ch.~\ref{ch:integer}\\
Positive feedback & A collision forces a positive elementary signed cycle. Cutting all such cycles gives a valid counting bound. & Ch.~\ref{ch:signed}\\
Graph obstruction & The universal subpolynomial-feedback criterion is conditionally incompatible with an explicit Bateman--Horn instance. Not a disproof of Erd\H{o}s 1060. & Ch.~\ref{ch:bh}\\
Genuine exchanges & Exact decomposition into disjoint minimal arithmetic exchanges; valid packing and weighted bounds. Uniform exchange count unproved. & Ch.~\ref{ch:exchanges}\\
\bottomrule
\end{longtable}
\normalsize

\section{What would actually complete the proof}
For fixed $E\ge1$, let
\[
 f_E(M)=\#\{k:h(k)=M,\ v_p(k)\le E\ \forall p\},\qquad
 F_E(X)=\max_{1\le M\le X}\max\{1,f_E(M)\}.
\]
The proved high-exponent reduction shows that it suffices to establish
\begin{equation}\label{eq:BE-front}
 \boxed{\log F_E(X)=o_E(\log X/\log\log X)
 \quad\text{for every fixed }E.}
\end{equation}
Conversely the full conjecture implies these estimates, after absorbing finitely
many small targets. Thus the all-cap statement is an equivalent reformulation,
not an independent theorem that has already solved most of the difficulty.
A proof only for $E=2$, or only under bounded small-prime target valuations, does
not suffice. Chapter~\ref{ch:reductions} contains the exact reduction and the
necessary order of quantifiers.

\section{Notation and conventions}
\begin{longtable}{@{}p{0.24\textwidth}p{0.7\textwidth}@{}}
\toprule Symbol & Meaning\\\midrule\endhead
$\alpha_p$ & $v_p(N)$, the target exponent, not necessarily an input exponent.\\
$e_p$ & $v_p(k)$, an input exponent; $e_p+1$ is its exponent index.\\
$\omega(n),\Omega(n)$ & Number of distinct prime factors; number with multiplicity.\\
$P^+(n),P^-(n)$ & Largest and smallest prime factor, with $P^-(1)=\infty$ when roughness is used.\\
$\mathcal P(N)$ & $\prod_{p\mid N}p/(p-1)$, which is $O(\log\log N)$.\\
$t_{\ell'}$ & $t/\ell^{v_\ell(t)}$, the $\ell$-free part.\\
$D_p$ & A domain of local input exponents, always required to preserve every genuine solution.\\
$Y$-rough & No prime factor at most $Y$; $1$ is included. A rough input and a rough divisor sum are different assumptions.\\
Compatible & At each base, every pair of exponent indices in the family is divisibility-comparable.\\
Unitary factor & $d\mid k$ with $\gcd(d,k/d)=1$. Only full equal-exponent blocks can automatically be canceled through $h$.\\
Positive cycle & An elementary directed signed cycle whose signs multiply to $+1$, not two traversals of a negative cycle.\\
Minimal exchange & An exact equality of products of local $h$-blocks, minimal for its specified changes. Not every literature use of ``primitive collision.''\\
\bottomrule
\end{longtable}

\section{Source and experiment labels}
\textbf{Working proof:} a complete argument preserved from this investigation,
with its hypotheses stated. \textbf{External theorem:} a result attributed to a
specific primary source. \textbf{Finite certificate:} an exact computation with a
specified scope and a checkable certificate. \textbf{Conditional:} a conclusion
requiring an unproved hypothesis such as Bateman--Horn. \textbf{Heuristic or open
route:} a proposal with no claimed proof.

A resource-limited search without a witness is inconclusive. A census through
input $K$ is not a census of every encountered larger target: it captures complete
fibers for targets at most $K^2$, since $h(k)>k^2$ for $k>1$. A finite prime-factor
range permits very large inputs but remains finite in allowed primes. The
experimental tables identify which kind of bound was used.
